\section{General linear path stability}
\label{sec:general-linear-path-stability}

Fable 5.2 proposed the localization and shortest-path layer argument in its November 10 research notes.
The Codex root and path-stability agents independently audited that argument.
The proof below includes the minor padding correction found in both audits.
It uses Theorem~\ref{thm:research-marked-size} to supply the required forest estimate.
Thus the conclusion is a mathematical theorem, rather than a conditional reduction.
It has not been formalized in Lean.

All logarithms in this section have base two, except $\ln$.
Put $\Phi(x)=x/\log^2x$ for $x\ge8$.
This function is increasing on its domain.

\begin{theorem}[General linear stability]\label{thm:general-linear-path-stability}
For every $0<\varepsilon\le1/2$, there is a constant $c_\varepsilon>0$ with the following property.
Let $N$ be a connected simple graph of rank $b$, with distinct terminals $u,v$.
If
\[
 P_{uv}(N)\ge(1/2+\varepsilon)2^b,
\]
then every $u$--$v$ path has length at least $c_\varepsilon b$.
In particular, $\operatorname{dist}_N(u,v)\ge c_\varepsilon b$.

An explicit expression in terms of the forest threshold is available.
Set $\delta=\varepsilon/2$, and let $R_\delta$ be an integer as in Lemma~\ref{lem:general-fe2} below.
Put
\[
 M=2R_\delta,\qquad
 \Theta=\max\left\{2^{36},\frac{(M+1)M2^{M+1}}{\delta}\right\}.
\]
One may take $c_\varepsilon=(\Theta+1)^{-2}$.
Every terminal path also satisfies the stronger inequality
\begin{equation}\label{eq:general-linear-strong}
 \ell>\frac{2b}{\Theta+1}-(\Theta+1).
\end{equation}
\end{theorem}

The formula determines a positive constant once the forest threshold is fixed.
The present proof does not give a numerical value for that threshold.

\subsection{The forest input and its localization}

\begin{lemma}\label{lem:general-fe2}
Fix $0<\delta<1/2$.
There is an integer $R_\delta\ge2^{100}$ such that the following statement holds.
If $G$ is connected and simple, $\rho=\beta(G)\ge R_\delta$, and $W$ is one cycle with
$|E(W)|\le\Phi(\rho)$, then
\[
 p(G,W)\le1/2+\delta.
\]
\end{lemma}

\begin{proof}
Use $a_\delta$ from Theorem~\ref{thm:research-marked-size}.
Choose $R_\delta$ above that theorem's rank threshold and so that
$\log R_\delta\ge\max\{100,1/a_\delta\}$.
For $\rho\ge R_\delta$,
\[
 \Phi(\rho)=\frac{\rho}{\log^2\rho}
 \le a_\delta\frac{\rho}{\log\rho}.
\]
The single-component condition also holds, since $1\le\log\rho/100$.
Theorem~\ref{thm:research-marked-size} applies.
\end{proof}

We also use the manuscript's small-cut estimate in the following form.
If $q$ disjoint cyclic regions avoid $V(W)$ and each has boundary at most an integer $D\ge1$, then
\begin{equation}\label{eq:general-small-cut-input}
 p(G,W)\le\frac12+\frac{(D+1)2^{D-1}}{q}.
\end{equation}
This statement permits loops and parallel edges.
It follows directly from the cut formula, the exceptional-partner bound, and the second-moment estimate.

\begin{lemma}[Localization]\label{lem:general-localization}
Let $G$ be a connected multigraph, and let $R$ be a nonempty connected proper vertex set.
Suppose that $G-R$ is connected, and write $d=|\delta_G R|\ge1$.
For each boundary edge $f$, with endpoint $z_f\in R$, add a distinct new vertex $t_f$ and an edge $z_ft_f$.
Add a cycle $W_R$ through the vertices $t_f$, padding it with new vertices if $d<3$.
Together with $G[R]$, this gives a graph $G_R$ satisfying
\[
 |E(W_R)|=\max\{d,3\},\qquad
 \beta(G_R)=\beta(G[R])+d.
\]
For uniform even subgraphs,
\begin{equation}\label{eq:general-localization-law}
 \Pr\bigl(\xi_G\cap E(G[R])\text{ is a forest}\bigr)=p(G_R,W_R).
\end{equation}
If $R\cap V(W)=\varnothing$, then $p(G,W)\le p(G_R,W_R)$.
\end{lemma}

\begin{proof}
Restriction outside an edge set $T$ maps the cycle space onto the cycle space after contracting every component of $T$.
For surjectivity, prescribe an even edge set in the contraction.
Its boundary inside each component of $T$ has even size.
A join in a spanning tree of that component completes the prescribed edge set.
This is a surjective linear map and therefore preserves uniformity on its image.

Contract $E(G-R)$ in $G$.
The resulting graph consists of $G[R]$ and one vertex receiving its $d$ boundary edges.
Contract $E(W_R)$ in $G_R$.
The same graph results, so the restrictions to $E(G[R])$ have the same law.
The added spokes have distinct leaf endpoints outside $W_R$.
They cannot create a cycle outside $W_R$, which proves \eqref{eq:general-localization-law}.
The stated inequality follows because $E(G[R])$ avoids $E(W)$.

For $d\ge3$, the construction adds $d$ vertices and $2d$ edges, giving the rank formula.
For $d=1,2$, padding the cycle adds equally many further edges and vertices.
Thus the rank formula remains valid.
\end{proof}

\begin{lemma}[Connected complementary shore]\label{lem:general-shore-closure}
Let $G$ and $W$ be connected, and let $R$ be a connected vertex set avoiding $V(W)$.
Let $C_0$ be the component of $G-R$ containing $W$, and put $R^+=V(G)\setminus V(C_0)$.
Then $R^+$ is connected, contains $R$, avoids $V(W)$, and satisfies
$\delta_G R^+\subseteq\delta_G R$.
Its complement is connected.
\end{lemma}

\begin{proof}
Every component of $G-R$ other than $C_0$ has an edge to $R$, since $G$ is connected.
Their union with $R$ is therefore connected.
Edges between $R^+$ and $C_0$ must have their other endpoint in $R$.
This gives the boundary inclusion and the remaining conclusions.
\end{proof}

\subsection{A shortest path in a cubic expansion}

\begin{lemma}\label{lem:general-shortest-reduction}
Suppose $P_{uv}(N)>2^{b-1}$.
There is a connected loopless multigraph $\Gamma$, with distinct terminals $u',v'$ and a shortest terminal path $A$, such that
\begin{enumerate}
\item $\beta(\Gamma)=b$ and $P_{u'v'}(\Gamma)\ge P_{uv}(N)$;
\item every nonterminal vertex of $\Gamma$ has degree three, and both terminals have degree one or two;
\item the length $\ell_0$ of $A$ is at most the length of every terminal path in $N$;
\item $|V(\Gamma)|\ge2b$, and $U=V(\Gamma)\setminus V(A)$ has size at least $2b-\ell_0-1$.
\end{enumerate}
Every vertex of $A$ has at most one neighbor incidence outside $A$.
\end{lemma}

\begin{proof}
By Proposition~\ref{prop:research-offroute}, no cyclic block lies outside the terminal route.
Delete the trees outside that route.
Suppress every degree-two vertex other than $u,v$.
This preserves rank and path count, and it can only shorten terminal paths.
No loop arises, since a cyclic block on the route retains its two distinct route terminals.
Call the resulting kernel $N_1$.
Every nonterminal vertex of $N_1$ has degree at least three.

Choose a shortest terminal path $A_1$ in $N_1$.
At each internal vertex of $A_1$, retain one vertex incident with its two path edges.
Attach all other incident edges through a rooted trivalent tree in its third branch.
At each terminal, retain one vertex incident with the path edge and attach its other edges through one rooted tree.
A single remaining incident edge needs no internal tree vertex.
Replace every other vertex of degree at least four by a trivalent tree.
No embedding restriction is imposed on these trees.

These replacements preserve rank.
Every original terminal path has a unique simple lift, so path count cannot decrease.
The image $A$ of $A_1$ has exactly the same length $\ell_0$.
Contracting the replacement trees maps any new terminal path to a terminal walk in $N_1$ of no greater length.
Every such walk has length at least the terminal distance in $N_1$, namely $\ell_0$.
Thus $A$ remains shortest.
This argument uses a shortest path after suppression, not an arbitrary path in the original graph.

All nonterminal degrees are three, and both terminal degrees are at most two.
If $n=|V(\Gamma)|$, the degree sum gives
\[
 2(n+b-1)=3(n-2)+\deg(u')+\deg(v').
\]
Hence $n=2b+4-\deg(u')-\deg(v')\ge2b$.
Subtracting $|V(A)|=\ell_0+1$ proves the last assertion.
The construction also gives the asserted bound on incidences from $A$ to $U$.
\end{proof}

\subsection{Wide distance layers give a region with small boundary}

Write $A=w_0w_1\cdots w_{\ell_0}$.
Let $L_j$ be the distance-$j$ layer from $u'=w_0$ in $\Gamma$.
Since $A$ is shortest, $w_j\in L_j$.
Put $U_j=L_j\cap U$ and $m_j=|U_j|$, taking $m_j=0$ for nonexistent layers.
For integers $i\le i'$, put
\[
 R(i,i')=U_i\cup\cdots\cup U_{i'},\qquad
 m(i,i')=\sum_{j=i}^{i'}m_j.
\]
Every edge joins vertices in the same or consecutive distance layers.
At most one edge joins each marked vertex to $U$.
Consequently,
\begin{equation}\label{eq:general-window-boundary}
 |\delta_\Gamma R(i,i')|
 \le3m_{i-1}+3m_{i'+1}+(i'-i+3).
\end{equation}
The last term counts the possible marked vertices in layers $i-1,\ldots,i'+1$.

For a parameter $\Theta\ge2^{36}$, call $R\subseteq U$ good if
\[
 |R|\ge\Theta,\qquad |\delta_\Gamma R|\le\Phi(|R|)/4.
\]
A good set need not be connected.

\begin{lemma}[Growing a window]\label{lem:general-window}
If $m_j\ge\Theta$ for some $j$, then a good set exists.
\end{lemma}

\begin{proof}
Start with the window $[j,j]$.
While a neighboring layer has mass greater than $\Phi(m(i,i'))/48$, include that layer.
Repeat until neither neighboring layer satisfies the condition.
Only finitely many layers are nonempty, so the procedure stops.

Let $\mu$ be the final mass, and set $L=\log\mu\ge36$.
Every step multiplies the current mass by more than $1+1/(48L^2)$.
Thus the number of steps is at most
\[
 \frac{\ln\mu}{\ln(1+1/(48L^2))}
 \le96L^2\ln\mu<67L^3.
\]
At termination, each neighboring layer has mass at most $\Phi(\mu)/48$.
Equation~\eqref{eq:general-window-boundary} therefore gives
\[
 |\delta_\Gamma R(i,i')|
 \le\Phi(\mu)/8+67L^3+3\le\Phi(\mu)/4.
\]
The final inequality follows from
$2^L\ge8L^2(67L^3+3)$ for $L\ge36$.
It holds at $L=36$, and the ratio of its left side to its right side increases thereafter.
Since $\mu\ge\Theta$, the final window is good.
\end{proof}

\begin{lemma}[No good set forces a long path]\label{lem:general-layer-count}
If no good subset of $U$ exists, then
\[
 \ell_0>\frac{2b}{\Theta+1}-(\Theta+1).
\]
\end{lemma}

\begin{proof}
Lemma~\ref{lem:general-window} gives $m_j<\Theta$ in every layer.
Let $R$ be the union of the layers $U_j$ with $j\ge\ell_0+2$.
No vertex of $A$ has a neighbor in $R$.
Therefore $|\delta_\Gamma R|\le3m_{\ell_0+1}<3\Theta$.
If $|R|\ge\Theta^2$, monotonicity of $\Phi$ and $\Phi(\Theta^2)/4\ge3\Theta$ would make $R$ good.
Hence $|R|<\Theta^2$.
Counting the remaining layers gives
\[
 2b-\ell_0-1\le|U|<(\ell_0+2)\Theta+\Theta^2.
\]
Rearranging proves the assertion.
\end{proof}

\subsection{A good set contradicts excess path density}

Add a new vertex $x$ and edges $u'x,xv'$ to $\Gamma$.
Call the resulting graph $H$, and mark the cycle $W=A\cup\{u'x,xv'\}$.
The proof of Lemma~\ref{lem:research-path-injection} also applies to multigraphs, giving
\begin{equation}\label{eq:general-path-forest}
 P_{u'v'}(\Gamma)\le p(H,W)2^b.
\end{equation}

\begin{lemma}[The good-set dichotomy]\label{lem:general-good-set}
Use $\delta,R_\delta,M,\Theta$ from Theorem~\ref{thm:general-linear-path-stability}.
If a good subset of $U$ exists, then $p(H,W)\le1/2+\delta$.
\end{lemma}

\begin{proof}
Let $R$ be good, with $m=|R|$ and $d=|\delta_\Gamma R|\le m/(4\log^2m)$.
Let $R_s$ be the vertex sets of the components of $\Gamma[R]$.
Their boundaries are disjoint and have total size $d$.
Every vertex in $R$ is cubic in both $\Gamma$ and $H$, so
\[
 \beta(\Gamma[R_s])=\frac{|R_s|-|\delta_\Gamma R_s|}{2}+1.
\]
Separate the components into four classes:
\begin{enumerate}
\item $|R_s|\ge M$ and $|\delta_\Gamma R_s|\le\Phi(|R_s|)/2$;
\item $|R_s|\ge M$ and $|\delta_\Gamma R_s|>\Phi(|R_s|)/2$;
\item $|R_s|<M$ and $\Gamma[R_s]$ is a tree;
\item $|R_s|<M$ and $\Gamma[R_s]$ contains a cycle.
\end{enumerate}

Class 2 has total size less than $2d\log^2m\le m/2$.
Indeed, each of its components satisfies
$|R_s|<2|\delta_\Gamma R_s|\log^2|R_s|\le2|\delta_\Gamma R_s|\log^2m$.
A tree component satisfies $|\delta_\Gamma R_s|=|R_s|+2$.
Thus class 3 has total size less than $d\le m/4$.
A cyclic component in class 4 has boundary at most its size, hence less than $M$.

Suppose first that class 1 is empty.
Class 4 then has total size greater than $m/4$, so it contains more than $m/(4M)$ components.
These are disjoint cyclic regions avoiding $W$, with boundaries less than $M$.
By \eqref{eq:general-small-cut-input},
\[
 p(H,W)\le\frac12+
 \frac{(M+1)2^{M-1}\,4M}{m}
 \le\frac12+\delta.
\]
The last inequality uses $m\ge\Theta$.

Suppose instead that class 1 contains $R_s$.
Apply Lemma~\ref{lem:general-shore-closure} in $H$ to obtain $R^+\supseteq R_s$ with connected complement.
It avoids $V(W)$, so it is contained in $U$ and all its vertices are cubic.
Write $d^+=|\delta_H R^+|$ and
\[
 \rho=\beta(H[R^+])+d^+
 =\frac{|R^+|+d^+}{2}+1
 \ge|R_s|/2\ge R_\delta.
\]
Since $\delta_HR_s=\delta_\Gamma R_s$, the boundary comparison gives
\[
 d^+\le|\delta_\Gamma R_s|
 \le\Phi(|R_s|)/2
 \le\Phi(|R_s|/2)
 \le\Phi(\rho).
\]
The middle inequality uses $\log(s/2)=\log s-1$.

Localize at $R^+$ using Lemma~\ref{lem:general-localization}.
If $d^+\ge3$, its marked cycle has at most $\Phi(\rho)$ edges by the last display.
If $d^+=1$ or $2$, its marked cycle has three edges, which still meets that bound since $\rho\ge2^{100}$.
Subdivide any parallel interior edges until the localized graph is simple.
Such subdivisions preserve rank, marked-cycle length, and outside-forest probability.
Lemma~\ref{lem:general-fe2} now yields
\[
 p(H,W)\le p(H_{R^+},W_{R^+})\le1/2+\delta.
\]
This completes both cases.
\end{proof}

\begin{proof}[Proof of Theorem~\ref{thm:general-linear-path-stability}]
Use Lemma~\ref{lem:general-shortest-reduction}.
The density assumption and $\delta=\varepsilon/2$ give
\[
 p(H,W)\ge\frac{P_{u'v'}(\Gamma)}{2^b}
 \ge\frac12+\varepsilon>\frac12+\delta.
\]
Lemma~\ref{lem:general-good-set} excludes every good subset of $U$.
Lemma~\ref{lem:general-layer-count} proves \eqref{eq:general-linear-strong} for $\ell_0$.
Every original terminal path has length at least $\ell_0$, so the same inequality holds for its length $\ell$.

Put $Q=\Theta+1$.
If $b\ge Q^2$, then $\ell>2b/Q-Q\ge b/Q\ge b/Q^2$.
If $b<Q^2$, then $\ell\ge1>b/Q^2$.
Thus $c_\varepsilon=Q^{-2}$ works at every rank.
\end{proof}

\subsection{Structural consequences}

\begin{corollary}\label{cor:general-short-path-density}
For any sequence of connected graphs with ranks $b\to\infty$ and terminal distances $o(b)$,
\[
 P_{uv}(N)\le(1/2+o(1))2^b.
\]
If $P_{uv}(N)\ge(1-\alpha)2^b$ for a fixed $0\le\alpha<1/2$, every terminal path has length
$\Omega_\alpha(b)$.
\end{corollary}

\begin{proof}
A fixed positive density gap in the first assertion would contradict Theorem~\ref{thm:general-linear-path-stability}.
For the second assertion, use $\varepsilon=1/2-\alpha$.
\end{proof}

\begin{corollary}[Bounded width of the terminal kernel]\label{cor:general-bounded-pathwidth}
Under the hypotheses of Theorem~\ref{thm:general-linear-path-stability}, remove all blocks outside the terminal route and suppress nonterminal degree-two vertices.
The resulting terminal kernel has pathwidth at most $2\lceil\Theta\rceil+1$.
The same bound holds for the cubic expansion in Lemma~\ref{lem:general-shortest-reduction}.
\end{corollary}

\begin{proof}
The proof excludes good sets, so Lemma~\ref{lem:general-window} gives $|U_j|<\Theta$ in every distance layer.
At most one marked path vertex lies in each layer.
Therefore each full layer has at most $\lceil\Theta\rceil+1$ vertices.
The bags $L_j\cup L_{j+1}$ form a path decomposition, since every edge joins equal or consecutive layers.
Their sizes give the stated pathwidth bound.
The terminal kernel is obtained from the cubic expansion by contracting its replacement trees.
Pathwidth does not increase under taking minors.
\end{proof}

The restriction to the terminal kernel matters.
Trees attached outside the terminal route can have arbitrarily large pathwidth without changing rank or terminal path count.
